\section{Repaso de c\'alculo en una dimensi\'on}
\label{sec:repaso_calculo}

\subsection{Funciones}
\label{subsec:funciones}

\begin{definicion}[Rango]
  \label{def:rango}
  El \textbf{rango} de una funci\'on $f:X\rightarrow Y$, con
  $X,Y\subset\R$, es el conjunto de los elementos de $Y$ que son valores que
  toma $f$:
  \begin{equation}
    \label{eq:rango}
    \mathrm{Ran}(f)=\set{y\in Y\such y=f(x)\ \text{para alg\'un}\ x\in X}.
  \end{equation}
\end{definicion}

\begin{definicion}[Sobreyectiva]
  \label{def:sobreyectiva}
  $f:X\rightarrow Y$ es \textbf{sobre} o \textbf{sobreyectiva} si cada elemento
  de $Y$ es imagen de alg\'un $x\in X$, esto es, si $\mathrm{Ran}(f)=Y$.
\end{definicion}

\begin{definicion}[Inyectiva]
  \label{def:inyectiva}
  $f:X\rightarrow Y$ es \textbf{uno a uno} o \textbf{inyectiva} si dos
  elementos distintos de $X$ tienen im\'agenes distintas en $Y$; es decir, si
  para cualesquiera $x_1,x_2\in X$ con $x_1\neq x_2$ se cumple
  $f(x_1)\neq f(x_2)$.
\end{definicion}

\subsection{L\'imites}
\label{subsec:limites_1d}

\begin{definicion}[L\'imite]
  \label{def:limite_1d}
  La funci\'on $f$ tiende al \textbf{l\'imite} $\ell$ en $x_0$ si
  $\forall\ep>0$ existe $\de_\ep>0$ tal que para todo $x$ con
  $0<\abs{x-x_0}<\de_\ep$ se cumple $\abs{f(x)-\ell}<\ep$. Simb\'olicamente,
  \begin{equation}
    \label{eq:limite_1d}
    \lim_{x\rightarrow x_0}f(x)=\ell .
  \end{equation}
\end{definicion}

\begin{ejemplo}
  \label{ej:limite_sandwich}
  Calculemos
  \begin{equation}
    \label{eq:limite_ejemplo}
    \lim_{x\rightarrow\infty}\frac{x+\sin^3x}{5x+6}=\frac{1}{5}.
  \end{equation}
  Como $-1\le\sin^3x\le 1$, la funci\'on queda encajonada,
  \begin{equation}
    \label{eq:sandwich}
    \frac{x-1}{5x+6}\le f(x)\le\frac{x+1}{5x+6}.
  \end{equation}
  Tomando $\xii=1/x$,
  \begin{equation}
    \label{eq:cambio_variable}
    \frac{x-1}{5x+6}
      =\frac{1-\frac{1}{x}}{5+\frac{6}{x}}
      =\frac{1-\xii}{5+6\xii},
  \end{equation}
  lo que transforma el l\'imite en $\lim_{\xii\rightarrow 0}$. Entonces
  \begin{equation}
    \label{eq:estimacion}
    \abs{\frac{1-\xii}{5+6\xii}-\frac{1}{5}}
      =\abs{\frac{-11\xii}{25+30\xii}}<\abs{\xii}<\de_\ep
  \end{equation}
  para $\xii$ suficientemente peque\~no, y basta tomar $\de_\ep=\ep$.
\end{ejemplo}

\subsection{Funciones continuas}
\label{subsec:continuidad_1d}

Para una funci\'on cualquiera no necesariamente se cumple
$\lim_{x\rightarrow x_0}f(x)=f(x_0)$.

\begin{definicion}[Continuidad]
  \label{def:continuidad_1d}
  $f$ es \textbf{continua} en $x_0$ si
  \begin{equation}
    \label{eq:continuidad_1d}
    \lim_{x\rightarrow x_0}f(x)=f(x_0),
  \end{equation}
  y es continua en un dominio $D$ si esto vale para todo $x_0\in D$.
\end{definicion}

\begin{teorema}
  \label{teo:algebra_continuas}
  Si $f$ y $g$ son continuas en $x_0$, entonces:
  \begin{enumerate}
    \item $f+g$ es continua en $x_0$;
    \item $fg$ es continua en $x_0$;
    \item si $g(x_0)\neq 0$, entonces $1/g$ es continua en $x_0$.
  \end{enumerate}
\end{teorema}

\begin{teorema}[Valor intermedio]
  \label{teo:valor_intermedio}
  Si $f$ es continua en $[a,b]$ y $f(a)<0<f(b)$, entonces existe alg\'un
  $x_0\in[a,b]$ tal que $f(x_0)=0$.
\end{teorema}

\begin{teorema}[Acotamiento]
  \label{teo:acotamiento}
  Si $f$ es continua en $[a,b]$, entonces $f$ est\'a acotada superiormente en
  $[a,b]$; es decir, existe $N\in\R$ tal que $f(x)\le N$ para todo
  $x\in[a,b]$.
\end{teorema}

\begin{teorema}[Valor extremo]
  \label{teo:valor_extremo}
  Si $f$ es continua en $[a,b]$, entonces existe alg\'un $y\in[a,b]$ tal que
  $f(y)\ge f(x)$ para todo $x\in[a,b]$.
\end{teorema}

\subsection{Derivaci\'on}
\label{subsec:derivacion_1d}

La pendiente de la recta secante que une los puntos de abscisas $x$ y $x+h$ es
$\qty{f(x+h)-f(x)}/h$. La pendiente de la recta \emph{tangente} en
$\qty{x,f(x)}$ se obtiene en el l\'imite.

\begin{definicion}[Derivada]
  \label{def:derivada_1d}
  $f$ es \textbf{derivable} en $x_0$ si existe
  \begin{equation}
    \label{eq:derivada_1d}
    \dv{}{x}f(x)=\lim_{h\rightarrow 0}\frac{f(x+h)-f(x)}{h}.
  \end{equation}
\end{definicion}

\begin{teorema}
  \label{teo:derivable_continua}
  Si $f$ es derivable en $x_0$, entonces $f$ es continua en $x_0$.
\end{teorema}

\begin{demostracion}
  Basta escribir la diferencia como el cociente incremental multiplicado por
  $h$:
  \begin{equation}
    \label{eq:demo_continuidad}
    \begin{aligned}
      \lim_{h\rightarrow 0}\qty{f(x+h)-f(x)}
        &=\lim_{h\rightarrow 0}\frac{f(x+h)-f(x)}{h}\,h\\
        &=\qty{\dv{}{x}f(x)}\lim_{h\rightarrow 0}h=0 ,
    \end{aligned}
  \end{equation}
  puesto que el primer factor existe y es finito por hip\'otesis. Por lo tanto
  $\lim_{h\rightarrow 0}f(x+h)=f(x)$ en $x_0$.
\end{demostracion}

\subsubsection{Consecuencias}
\label{subsubsec:reglas_derivacion}

Sean $f$ y $g$ derivables en $x\in(a,b)$. Entonces
\begin{align}
  \label{eq:regla_suma}
  \dv{}{x}(f+g)&=\dv{}{x}f+\dv{}{x}g,\\
  \label{eq:regla_producto}
  \dv{}{x}(fg)&=\qty{\dv{}{x}f}g+\qty{\dv{}{x}g}f,
\end{align}
y si $g(x_0)\neq 0$,
\begin{equation}
  \label{eq:regla_reciproco}
  \dv{}{x}\qty{\frac{1}{g}}=-\qty{\dv{}{x}g}\frac{1}{g^2}.
\end{equation}

\begin{demostracion}[Demostraci\'on de \eqref{eq:regla_suma}]
  \begin{equation}
    \label{eq:demo_suma}
    \begin{aligned}
      \dv{}{x}(f+g)
        &=\lim_{h\rightarrow 0}\frac{f(x+h)+g(x+h)-f(x)-g(x)}{h}\\
        &=\lim_{h\rightarrow 0}\frac{f(x+h)-f(x)}{h}\\
        &\quad+\lim_{h\rightarrow 0}\frac{g(x+h)-g(x)}{h}\\
        &=\dv{}{x}f+\dv{}{x}g .
    \end{aligned}
  \end{equation}
\end{demostracion}

\begin{demostracion}[Demostraci\'on de \eqref{eq:regla_producto}]
  Se suma y se resta el t\'ermino cruzado $f(x+h)g(x)$:
  \begin{equation}
    \label{eq:demo_producto}
    \small
    \begin{aligned}
      \dv{}{x}(fg)
        &=\lim_{h\rightarrow 0}\frac{f(x+h)g(x+h)-f(x)g(x)}{h}\\
        &=\lim_{h\rightarrow 0}\frac{1}{h}
          \big(f(x+h)g(x+h)-f(x+h)g(x)\\
        &\qquad\qquad+f(x+h)g(x)-f(x)g(x)\big)\\
        &=\lim_{h\rightarrow 0}f(x+h)\frac{g(x+h)-g(x)}{h}\\
        &\quad+\lim_{h\rightarrow 0}\frac{f(x+h)-f(x)}{h}g(x)\\
        &=f(x)\dv{}{x}g(x)+\qty{\dv{}{x}f(x)}g(x).
    \end{aligned}
  \end{equation}
\end{demostracion}

\begin{demostracion}[Demostraci\'on de \eqref{eq:regla_reciproco}]
  \begin{equation}
    \label{eq:demo_reciproco}
    \begin{aligned}
      \dv{}{x}\qty{\frac{1}{g}}
        &=\lim_{h\rightarrow 0}\frac{1}{h}
          \qty{\frac{1}{g(x+h)}-\frac{1}{g(x)}}\\
        &=\lim_{h\rightarrow 0}\frac{g(x)-g(x+h)}{h\,g(x)g(x+h)}\\
        &=-\lim_{h\rightarrow 0}\frac{g(x+h)-g(x)}{h}
          \lim_{h\rightarrow 0}\frac{1}{g(x)g(x+h)}\\
        &=-\qty{\dv{}{x}g(x)}\frac{1}{g(x)^2}.
    \end{aligned}
  \end{equation}
\end{demostracion}

\subsubsection{Regla de la cadena}
\label{subsubsec:regla_cadena_1d}

Si $f$ y $g$ son derivables en $x_0$ y $k(x)=(f\circ g)(x)=f(g(x))$, entonces
\begin{equation}
  \label{eq:regla_cadena_1d}
  \dv{}{x}f(g(x))=\qty{\eval{\dv{f}{u}}_{u=g(x)}}\dv{}{x}g(x).
\end{equation}

\begin{ejemplo}[Derivada de $x^n$]
  \label{ej:derivada_potencia}
  Sea $f(x)=x^n$ con $n\in\N$. Por el binomio de Newton,
  \begin{equation}
    \label{eq:derivada_potencia}
    \begin{aligned}
      \dv{}{x}f(x)
        &=\lim_{h\rightarrow 0}\frac{(x+h)^n-x^n}{h}\\
        &=\lim_{h\rightarrow 0}\frac{1}{h}
          \qty{x^n+nx^{n-1}h+\dots+h^n-x^n}\\
        &=\lim_{h\rightarrow 0}\qty{nx^{n-1}+\mathcal{O}(h)}
         =nx^{n-1}.
    \end{aligned}
  \end{equation}
\end{ejemplo}

\subsection{Teoremas del valor medio}
\label{subsec:valor_medio}

\begin{definicion}[Extremo local]
  \label{def:extremo_local}
  Sea $f$ una funci\'on y $A$ un conjunto contenido en su dominio. Un punto
  $x\in A$ es un \textbf{m\'aximo} (o \textbf{m\'inimo}) \textbf{local} de $f$
  sobre $A$ si existe $\de>0$ tal que $x$ es un m\'aximo (o m\'inimo) de $f$
  sobre $A\cap\qty{x-\tfrac{\de}{2},\,x+\tfrac{\de}{2}}$.
\end{definicion}

\begin{teorema}[Rolle]
  \label{teo:rolle}
  Si $f$ es continua en $[a,b]$, derivable en $(a,b)$ y $f(a)=f(b)$, entonces
  existe $x_0\in(a,b)$ tal que
  \begin{equation}
    \label{eq:rolle}
    \eval{\dv{}{x}f(x)}_{x_0}=0 .
  \end{equation}
\end{teorema}

\begin{teorema}[Valor medio]
  \label{teo:valor_medio}
  Si $f$ es continua en $[a,b]$ y derivable en $(a,b)$, entonces existe
  $x_0\in(a,b)$ tal que
  \begin{equation}
    \label{eq:valor_medio}
    \eval{\dv{}{x}f(x)}_{x_0}=\frac{f(b)-f(a)}{b-a}.
  \end{equation}
\end{teorema}
